\documentclass[10pt]{article}
\usepackage[margin=20mm]{geometry}
\usepackage{amsmath,amssymb,mathtools,fontspec,unicode-math,xurl,hyperref}
\IfFontExistsTF{DejaVu Serif}{\setmainfont{DejaVu Serif}}{\setmainfont{Latin Modern Roman}}
\setmathfont{latinmodern-math.otf}
\setlength{\parindent}{0pt}\setlength{\parskip}{6pt}\sloppy
\title{Variational Domains and Magnetic Spectral Asymptotics}
\author{AI-assisted technical review}
\date{}
\begin{document}
\maketitle
Abstract: The frozen author manuscript arXiv:2306.05008v2 is examined. One Major defect and seven Minor corrections are established. The Major defect is an empty affine minimization class at the polynomial comparison step used for the explicit two-pole coefficients. A cutoff repairs the missing variational definition; the actual Fourier coefficients and principal eigenvalue expansion are not refuted. The remaining corrections concern scalar branches, a zero-order profile, boundary estimates, minimization, rescaled tests, Fourier normalization, and complex phase. This assessment concerns only the frozen author version. The scientific decision remains pending.\par
\section*{Source identity and scope}
Frozen identifier: M25-0259. Authors: Veronica Felli, Benedetta Noris, Roberto Ognibene, and Giovanni Siclari. Article: Quantitative spectral stability for Aharonov–Bohm operators with many coalescing poles. The frozen DOI is\par
https://doi.org/10.4171/JEMS/1632.\par
The inspected source has 44 pages and 664 840 bytes, with SHA256\par
`06a050eb0c244a7f8d569276c8c08800bdbd4a75e43410dbc33b58d7985e6ceb`.\par
Its arXiv version is v2, submitted on 9 December 2023 at 22:04:09 UTC. The first-page stamp says 9 December 2023, while the internal revision date is 7 December 2023; these different date fields are retained. The publisher records submission on 3 June 2023, acceptance on 18 April 2024, and online publication on 21 March 2025. Its subscription PDF was not acquired. All locations below refer to the frozen author PDF, not to an assumed pagination or correction status of the journal article. The frozen 2025 membership, identifier, and DOI are unchanged.\par
The setting includes a bounded simply connected planar domain containing the collision point, half-integer magnetic circulations, fixed distinct ray directions, and a simple limiting eigenvalue for the principal asymptotic theorems. The reduction uses real functions with opposite traces across the cracks. In particular, the two-pole closed formulas require the equal-distance, opposite-pole configuration (2.36). These hypotheses are retained throughout this review.\par
\section*{Confirmed corrections}
\subsection*{P1. The explicit real scalar branch is not smooth on the asserted domain (Minor)}
Location. Equation (2.8), p. 7; the angle choices and equation (3.5), p. 14.\par
Witness. Take \(\Omega=D_1\), \(k_1=0\), \(k_2=2\), equal radii \(r=1/8\), and \[(\alpha^1,\alpha^2,\alpha^3,\alpha^4)
       =(0,-\pi/2,\pi,\pi/2).\] This satisfies the stated indexing and opposite-direction requirements. The two horizontal angles, with indices \(1,3\), have the same positive sign in (3.5). At \(x=(-2r\varepsilon,0)\), which lies outside the finite horizontal segment and outside the vertical segment, both angle functions jump by \(2\pi\) in the same direction. The vertical angles are smooth there. Consequently \(\Theta_\varepsilon\) has a jump of magnitude \(2\pi\) at a point of \(\mathbb R^2\setminus\Gamma_\varepsilon\) and does not satisfy the first assertion of (2.8).\par
Counterevidence and effect. The jump disappears in \(e^{i\Theta_\varepsilon}\). On the actual crack segments the exponential has the required sign change, and its local logarithmic derivative gives the prescribed magnetic potential. Thus this example does not invalidate the circle-valued gauge or the spectral equivalence. There is also no separate-pair period obstruction: the segments form a connected cross, so a loop avoiding the cross cannot isolate one entire pair while excluding the other.\par
Repair boundary. Formulate the gauge directly as a circle-valued function with the appropriate local derivatives, or correct the scalar branch by multiples of \(2\pi\). The literal formula should not simultaneously be retained as a globally smooth real function on the larger claimed domain.\par
\subsection*{P2. The even-case profile omits the nonvanishing case (Minor)}
Location. Proposition 3.7 and (3.18), p. 17.\par
Witness. Set \(k_1=0\), \(k_2=1\), \(\Omega=D_1\), and take any allowed small opposite pair of poles. The limiting problem is the ordinary Dirichlet Laplacian. Its positive normalized first eigenfunction is \[v(x)=cJ_0(j_{0,1}|x|),\qquad c>0,\] and the first eigenvalue is simple. Since \(v(0)=c\ne0\), its vanishing order is \(m=0\). Formula (3.18) then gives the zero function on the right, whereas the left converges to \(c\). The interval containing \(\alpha_0\) also contains division by zero. Choosing \(m\ge1\) cannot rescue the assertion, since \(\varepsilon^{-m}v(\varepsilon\,\cdot)\) diverges.\par
Counterevidence, effect, and repair. The manuscript explicitly treats \(m=0\) separately in (3.19), Propositions 4.1 and 4.3, and the discussion preceding (2.38). Thus the nonvanishing case was intended and is not excluded by a hidden assumption. State its leading profile as a nonzero constant times the gauge sign function; reserve the sine profile and \(\alpha_0\) for \(m\ge1\). The later logarithmic energy bound does not require the erroneous zero sine profile. No claimed nonzero-order two-pole coefficient is refuted by this example.\par
\subsection*{P3. Interior analyticity does not give a global gradient bound on the stated domains (Minor)}
Location. The universal spatial quantifier in (3.19) and its proof, p. 17. Compare the corresponding local singular expansion on pp. 15–16.\par
Witness. Put \(b=(0,-1/2)\) and \[\Omega=\{b+\rho(\cos\theta,\sin\theta):
             0<\rho<1,\ 0<\theta<3\pi/2\}.\] This bounded simply connected domain contains \(D_{1/4}(0)\). Choose \(k_1=0\), \(k_2=1\), and poles initially at \((1/8,0)\) and \((-1/8,0)\). With \(\nu=2/3\), its positive normalized Dirichlet ground state is \[v\bigl(b+\rho(\cos\theta,\sin\theta)\bigr)
   =cJ_\nu(j_{\nu,1}\rho)\sin(\nu\theta).\] Separation of variables and positivity identify the simple first eigenvalue. At the collision point \(0\) it is positive (\(\rho=1/2\), \(\theta=\pi/2\)), so \(m=0\). The expansion \(J_\nu(j_{\nu,1}\rho)
 = (j_{\nu,1}/2)^\nu\rho^\nu/\Gamma(\nu+1)+O(\rho^{\nu+2})\) shows that \(|\nabla v|\) is unbounded like \(\rho^{-1/3}\) near the reentrant corner \(b\). This contradicts the constant gradient bound asserted for every point of \(\Omega\) in (3.19).\par
Counterevidence and effect. The needed bounds concern the small crack segments near \(0\), away from \(\partial\Omega\). On a fixed interior disk they follow from the displayed local expansions. The counterexample refutes the even-case global bound; it is not separately presented as an odd-case counterexample to (3.12).\par
Repair boundary. Localize the estimates to a fixed neighbourhood of the collision point, as sufficient for the asymptotic argument. A global estimate would require additional boundary regularity and an appropriate boundary estimate; interior analyticity alone supplies neither.\par
\subsection*{P4. Minimization on a half-line is replaced by minimization on the whole line (Minor)}
Location. The last displayed identity in the proof of Proposition 6.5, p. 26, after (6.12).\par
For fixed nonzero \(w\), write the displayed quadratic as \(At^2/2+Bt+C\), with \(A=\int|\nabla w|^2>0\). The asserted value \(C-B^2/(2A)\) is the minimum over \(t\in\mathbb R\), but the stated domain is \(t\ge0\). For \(B>0\) the actual infimum is \(C\).\par
Witness in the stated spaces. Choose a smooth bump \(w\) supported in a ball inside \(D_\varepsilon\) and away from all cracks, where \(v_0\) has a fixed nonzero sign, and give \(w\) the same sign. Then \(w\in\widetilde{\mathcal H}_\varepsilon\), \(\eta_\varepsilon=1\) on its support, and \(L_\varepsilon(w)=0\). The interior eigenvalue equation gives \[B=\int\nabla w\cdot\nabla v_0
       =\lambda_0\int wv_0>0 .\] This is a counterexample with an admissible direction, not just an unrelated scalar quadratic.\par
Counterevidence, effect, and repair. The outer optimization includes both \(w\) and \(-w\). Consequently the statement (6.11) remains correct: let \(t\) range over \(\mathbb R\) in the proof, or retain the one-sided minimum and then optimize over both signs. The erroneous fixed-direction equality does not change the supremum in (6.11).\par
\subsection*{P5. Zero extension does not preserve the variational equation for every whole-plane test (Minor)}
Location. Equation (6.20), p. 27, and its use in the proof of Proposition 6.7, p. 28.\par
The rescaled potential is supported in \(\Omega/\varepsilon\). The original weak equation (4.5) permits tests with the corresponding zero outer trace. It does not imply (6.20) for every \(w\in\widetilde{\mathcal H}\) on the plane.\par
Witness. Use the explicit simple half-flux rectangular model constructed under P8, with \(k_1=1\), \(k_2=0\), and the pole moving on the positive axis. Its real \(v_0\) is positive on a short positive-axis interval. The affine trace condition \(T(V_\varepsilon)=2v_0\) there forces \(V_\varepsilon\ne0\). If (6.20) held for all the stated tests, then in particular the zero extension of \(\widetilde V_\varepsilon\) would be weakly harmonic against every \(C_c^\infty(\mathbb R^2\setminus\Gamma_1)\) test: the right side vanishes for these tests. The slit plane is connected. A compactly supported harmonic function there vanishes identically, by interior regularity and unique continuation from its exterior zero set. This contradicts the nonzero affine trace. Thus at least one permitted whole-plane test fails the asserted equality.\par
Counterevidence and effect. For each fixed compactly supported test, its support is inside \(\Omega/\varepsilon\) once \(\varepsilon\) is sufficiently small. Use those tests first, pass to the limit, and only then use the density from Proposition 6.3. The energy test \(\widetilde V_\varepsilon-\eta\widetilde V_{0,\varepsilon}\) also has valid outer trace for sufficiently small \(\varepsilon\). These observations provide the limiting variational equation and strong energy convergence actually needed for Theorem 2.2.\par
Repair boundary. Restrict the finite-\(\varepsilon\) test space before taking the limit. Also replace the complement of \(\Omega\) immediately after (6.18) by the complement of \(\Omega/\varepsilon\). A zero extension generally introduces an outer boundary flux; it is not a solution on the whole plane at fixed \(\varepsilon\).\par
\subsection*{P6. The polynomial comparison functions are outside the defined affine domain (Major)}
Location. The assumption \(h\in H_0^1(\Lambda)\) and (7.1), pp. 34–35; the uses \(h=P_m\) and \(h=T_m\) in Propositions 7.2, 7.5, and 7.6, pp. 35–40.\par
Witness. Take \(P_1(x)=x_1\) and a disk or one of the ellipses \(E_\varepsilon(L)\) containing the segment. If \(V\in\mathcal H_{\varepsilon,\Lambda}\) and \(V-P_1\in\widetilde{\mathcal H}_{\varepsilon,\Lambda}\) as literally required, both have zero outer trace. Subtraction forces the outer trace of \(P_1\) to vanish, which it does not. The admissible set is empty. The same issue occurs for a nonzero homogeneous Taylor polynomial. In addition, an arbitrary \(L>0\) ellipse need not contain the fixed disk \(D_R\) required in the initial generic definition.\par
Counterevidence and effect. The differential problem (7.2) uses \(h\) only through its value and normal derivative on the crack, and its outer boundary condition concerns \(V\). This is the well-posed local problem used in the subsequent Fourier calculation, including (7.11) on p. 36. It does not retroactively make the earlier affine class for the unmodified polynomial nonempty or redefine its minimum. Propositions 7.2, 7.5 and 7.6 use that undefined variational object in the derivation of the explicit two-pole asymptotics. The failed identification is therefore a Major defect in this substantive proof chain. No error in the actual Fourier coefficients is inferred from the empty-class observation.\par
Repair boundary. Choose a smooth compactly supported \(\chi\) equal to one in a neighbourhood of the segment and use \(\chi P_m\) or \(\chi T_m\) in the affine constraint. The line functional and the solution of (7.2) are unchanged. A common cutoff in the smaller comparison domain preserves the zero-extension inclusions used in the domain sandwich. For a fixed \(L>0\), containing the segment in the interior, rather than containing the originally fixed \(D_R\), suffices. These repairs give the same computed leading constants; the literal empty minimization problem must nevertheless be corrected.\par
\subsection*{P7. The named Fourier basis has a different normalization (Minor)}
Location. The final paragraph of p. 36, after the definitions of \(a_{j,\varepsilon}\) and \(b_{j,\varepsilon}\).\par
The manuscript calls \(a_{j,\varepsilon}(0)\) coefficients relative to the orthonormal basis containing \(\cos(j\eta)/\sqrt\pi\). Its integral definition instead uses the ordinary cosine coefficients \((1/\pi)\int f(\eta)\cos(j\eta)\,d\eta\). For \(m=1\), \(\ell_0=1\), the boundary function is \(\cos\eta\): the defined coefficient is \(a_{1,\varepsilon}(0)=1\), whereas its orthonormal coefficient is \(\sqrt\pi\).\par
Counterevidence, effect, and repair. Equations (7.14) and (7.15) consistently use the ordinary coefficients and the correct factors \(\pi\) and \(\pi/2\) in Parseval’s identity. There is no propagated error in (7.6). Correct the description of the basis, or rescale all coefficients consistently; the first option leaves the subsequent calculation unchanged.\par
\subsection*{P8. The arbitrary complex phase is not tied to the earlier real profile (Minor)}
Location. The fixed real \(v_0\) in (2.15), p. 8; the real profile (2.26), p. 9; Theorem 2.4 and (2.32)–(2.33), pp. 10–11; its proof, p. 34.\par
For a compatible choice \(u_0^*=e^{i\Theta_0}v_0\), let \(u_\varepsilon^*=e^{i\Theta_\varepsilon}v_\varepsilon\) have the required relative normalization. Replacing both by \(iu_0^*,iu_\varepsilon^*\) preserves their norms and the positive real inner product in (2.32). It multiplies the blow-up limit by \(i\), whereas the earlier fixed real \(\Psi_0\) and its real minimizer \(\widetilde V\) leave the right side of (2.33) unchanged. The profile \(\Psi_0-\widetilde V\) is nonzero: \(\Psi_0\) is homogeneous of positive degree and has infinite Dirichlet energy at infinity, while \(\widetilde V\) has finite Dirichlet energy. Therefore these two limits cannot agree.\par
A model satisfying simplicity. Let \(\Omega=(-1/4,1)\times(-1,1)\), with a half-flux pole at \(0\). In the negative-axis cut formulation, reflection in the horizontal axis splits the spectrum into two problems on \(Q=(-1/4,1)\times(0,1)\), with Dirichlet conditions on the outer sides and top. The even sector has Dirichlet bottom for \(x_1<0\) and Neumann bottom for \(x_1>0\) (DN); the odd sector has the reversed bottom conditions (ND). A DN test equal to \(\sin(\pi x_1)\cos(\pi x_2/2)\) on \(x_1>0\), and zero on \(x_1<0\), gives \(\lambda_{\rm DN}\le5\pi^2/4\). A Neumann cut at \(x_1=0\) in the ND problem gives separate rectangular first eigenvalues \(17\pi^2/4\) and \(5\pi^2/4\). Thus \(\lambda_{\rm ND}\ge5\pi^2/4\), with strict inequality: equality would make the left piece zero and the right proportional to \(\cos(\pi x_1/2)\sin(\pi x_2)\), contradicting matching \(H^1\) traces at the interface. The DN ground state is positive and simple, so the full half-flux ground eigenvalue is simple. Choose \(R=1/8\) and \(a^1=(1/16,0)\). Its even extension is positive along the short positive-axis segment, also supplying the model used in P5.\par
This example supplies a simple branch for every stated \(\varepsilon\in(0,1]\), not just a small subsequence. For a moving pole at \(\delta=\varepsilon/16\), the DN bottom Dirichlet interval has length \(1/4+\delta\). Reflection of the ND problem in the vertical midline gives a DN problem with Dirichlet length \(1-\delta>1/4+\delta\). Strict ground-state monotonicity under enlarging the Dirichlet interval gives \(\lambda_{\rm DN}(\delta)<\lambda_{\rm ND}(\delta)\). Strictness follows from uniqueness and the impossibility of simultaneous zero value and normal derivative on an open boundary interval for a nonzero eigenfunction. The positive DN ground states have positive mutual inner product. Thus the family satisfies the normalization and simplicity hypotheses without numerical assumptions.\par
Counterevidence, effect, and repair. Remark 3.5 permits a real eigenbasis, and Section 4 explicitly chooses compatible real phases. The proof on p. 34 also redefines \(v_0\) from \(u_0\); these are strong evidence of the intended convention, but a relative inner product alone does not make that \(v_0\) real. Specify \(u_0=e^{i\Theta_0}v_0\) with the previously fixed real \(v_0\), or complexify and transform both profiles with the same phase. The eigenvalue formula and the energy norm in (2.34) are invariant under this change. The correction concerns the phase of the stated vector limit, not a new eigenvalue exponent or a different physical eigenspace.\par
\section*{Checks of the principal arguments}
\subsection*{General spectral expansion}
The real variational problem has a strictly coercive quadratic part on its affine trace class. Compactness can be checked by zero extension into a fixed disk cut into finitely many sectors; this also avoids assuming an unstated smooth outer boundary. The logarithmic cutoff at the collision point preserves opposite traces, and the weak-domain convergence and density argument give \[\|V_\varepsilon\|_{L^2}=o(\|V_\varepsilon\|_{\mathcal H_\varepsilon}).\] The latter argument only tests against functions vanishing near the shrinking segments before using density.\par
The exact identity central to Section 5 is \[\lambda_0(V_\varepsilon,v_0)_{L^2}
      =2\mathcal E_\varepsilon-2L_\varepsilon(v_0).\] It follows by combining the crack boundary terms with the equation for \(\psi_\varepsilon=v_0-V_\varepsilon\). The complementary spectral gap is uniform for sufficiently small \(\varepsilon\), since the limiting eigenvalue is simple. The spectral residual gives \(|\lambda_\varepsilon-\lambda_0|=O(\|V_\varepsilon\|_{L^2})\) and \(\|\psi_\varepsilon-\Pi_\varepsilon\psi_\varepsilon\|_{\mathcal H_\varepsilon}
 =O(\|V_\varepsilon\|_{L^2})\). Normalization then gives \(\|\widehat v_\varepsilon-v_0\|_{L^2}=O(\|V_\varepsilon\|_{L^2})\). Substitution in (5.17) yields the stated expansion, with an error bounded by \(O(\|V_\varepsilon\|_{L^2}^2)\), hence \(o(\|V_\varepsilon\|_{\mathcal H_\varepsilon}^2)\). If \(V_\varepsilon=0\) at an exceptional parameter, the weak equation already makes \(v_0\) a perturbed eigenfunction; divisions by its norm in a contradiction argument are needed only on a subsequence with \(V_\varepsilon\ne0\).\par
\subsection*{Odd number of poles}
The annular Hardy inequality relies on an odd number of sign changes: a zero-gradient function would be constant on each sector and the cyclic trace equations would force every constant to be zero. Summing annuli, and a local trace estimate, control both the weighted tail and the interior \(L^2\) norm. This establishes a genuine Hilbert norm and the cutoff density used for the limiting problem. With the restricted testing explained under P5, weak convergence identifies the unique limit; admissible energy tests then give strong convergence. The scaling is \(\varepsilon^{-m/2}\) for eigenfunctions and \(\varepsilon^{-m}\) for energy. Here \(m\) counts nodal half-lines and the actual vanishing order is \(m/2\); the phrase vanishing order \(m\) in Theorem 2.2 should be read in light of (2.21).\par
For poles on nodal rays the trace datum is zero and the nonzero linear functional makes the minimum negative. For bisector directions the linear functional vanishes and the nonzero affine trace makes the minimum positive. The rotation argument changes the data on a fixed space; dominated convergence and coercivity give the required continuity. It gives a zero leading coefficient between these configurations, but supplies no next nonzero exponent in that case. The manuscript itself acknowledges this limit.\par
\subsection*{Two opposite poles: exact constants}
The comparison problem is understood with the corrected admissibility in P6. Put \(a=\varepsilon r_1\) and \(T=\operatorname{arsinh}(L/a)\). Separation in elliptic coordinates gives the exact energy \[\int |\nabla V|^2
 =\pi a^{2m}\left\{
 \frac{\ell_0^2c_0^2}{2T}
 +\ell_0^2\sum_{j=1}^m j c_j^2\coth(jT)
 +\ell_1^2\sum_{j=1}^m\frac{d_j^2}{j}\tanh(jT)
 \right\}.\] The line functional on the solution equals \[L(V)=-\pi a^{2m}\ell_1^2
                 \sum_{j=1}^m\frac{d_j^2}{j}\tanh(jT).\] The cosine modes have prescribed value at \(\xi=0\) and zero value at \(\xi=T\); the sine modes have prescribed derivative at \(\xi=0\) and zero value at \(\xi=T\). These facts verify the signs independently of an asymptotic transcription. The zero mode contributes only the displayed \(1/T\) term.\par
For \(j>0\) of the same parity as \(m\), \[c_j=2^{1-m}\binom m{(m-j)/2},\qquad d_j=\frac jm c_j,\] and the other coefficients vanish. The identity \[\frac{m-2h}{m}\binom mh^2
   =\binom{m-1}h^2-\binom{m-1}{h-1}^2\] telescopes for \(0\le h\le\lfloor(m-1)/2\rfloor\) and proves (7.16); the formula for \(d_j\) proves (7.17). Own exact rational computations for \(1\le m\le32\) also checked the transcription. Those finite computations supplement, rather than replace, this proof for every \(m\).\par
The ellipse sandwich uses no reflection symmetry of \(\Omega\). Taylor’s remainder gives \(\|V_\varepsilon-V_{\varepsilon,T_m,\Omega}\|^2=O(\varepsilon^{2m+1})\). In (7.27), \(p>2\) implies \(1/p'>1/2\), so the remaining exponent \(2m-1/2+1/p'\) is strictly larger than \(2m\). The nodal case has \(\ell_0=0\), \(\ell_1=\pm m\beta\); the bisector case has \(\ell_1=0\), \(\ell_0=\pm\beta\). Combining the exact energy with the factor \(2\) in Theorem 2.1 recovers both coefficients \[\mp\frac{m\pi\beta^2 r_1^{2m}}{4^{m-1}}
       \binom{m-1}{\lfloor(m-1)/2\rfloor}^{\!2}.\] No error in these leading coefficients was found.\par
\section*{Supporting checks}
The initial reviewer read the primary normalized text in full, 3 260 lines with zero remaining; the 508-line source manifest was also read completely. The current report was compiled with the existing portable Tectonic route and its final pages were actually viewed before completion. Rendering alone was not recorded as inspection.\par
Their proofs, source-page checks, draft limitations, and final receipts are retained. The initial root reviewer read their final notes, checked the decisive arguments, constructed the explicit simple-eigenvalue model, and authored this report. The external controller’s recorded Personal/Chat/GPT6/Pro observations are attributed separately. These historical native AI checks did not constitute independent acceptance or human review. The current focused independent acceptance and its read-only auxiliary contribution are separately attributed in the verification files.\par
Section 8 gives specified collinear extensions and explicitly leaves the mixed odd/even configuration with poles on both sides open. It is not a proof of every collinear configuration. The proof audit uses the stated spectral continuity and local singular-expansion dependencies; their entire external papers were not independently re-proved. The internal use of those dependencies was checked. The closed combinatorial identity above was re-proved directly.\par
A bounded search of arXiv, EMS Press, CVGMT, the author’s publication page, and the Politecnico repository found no separate paper-specific supplement or code. The repository reported no associated files. Search results are not proof of universal absence; lack of code is not a mathematical defect. The public author/PDF version was not promoted to journal-version equivalence.\par
The current assessment is one Major and seven Minor findings. P6 breaks the written variational definition used for a central explicit asymptotic calculation; its repairability does not lower the severity of that derivation failure. The remaining local corrections have the stated nonpropagation evidence, including the valid exponential gauge, compact limiting tests and correct Fourier weights. No unproved candidate is counted as a confirmed defect. The scientific decision remains pending; `human\_reviews=[]`. The independent acceptance is separately recorded against the current artifacts.\par
9 V. Felli, B. Noris, R. Ognibene, and G. Siclari, Quantitative spectral stability for Aharonov–Bohm operators with many coalescing poles, arXiv:2306.05008v2, 9 December 2023. https://arxiv.org/abs/2306.05008v2. EMS Press, article metadata and access record, https://ems.press/journals/jems/articles/14298681. CVGMT, accepted-paper entry, https://cvgmt.sns.it/paper/6076/. R. Ognibene, publication list, https://sites.google.com/view/ognibene/publications. Politecnico di Milano, repository record, https://re.public.polimi.it/handle/11311/1298642.\par
\end{document}
